\documentclass[11pt,a4paper]{article}

% ─── Packages ───
\usepackage[utf8]{inputenc}
\usepackage[T1]{fontenc}
\usepackage[english]{babel}
\usepackage[margin=2.5cm]{geometry}
\usepackage{hyperref}
\usepackage{booktabs}
\usepackage{multirow}
\usepackage{graphicx}
\usepackage{caption}
\usepackage{enumitem}
\usepackage{parskip}

\hypersetup{
    colorlinks=true,
    linkcolor=blue!70!black,
    urlcolor=blue!70!black,
    citecolor=blue!70!black
}

\title{\textbf{Mid-Phase Progress Report} \\ \large CodeEnhancer Model Comparison}
\author{
    \normalsize \textbf{Group Members:} Ali Ekin Özçetin, Ali Sezer Uysal, Mert Güner \\
    \normalsize \textbf{Course:} CENG 467 — Natural Language Processing \\
    \normalsize \href{https://github.com/Aliekinozcetin/CodeEnhancer}{github.com/Aliekinozcetin/CodeEnhancer}
}
\date{}

\begin{document}

\maketitle

% ════════════════════════════════════════════
\section{Project Title, Group Members, and Problem Definition}
% ════════════════════════════════════════════

\textbf{Project Title:} CodeEnhancer Model Comparison

\textbf{Group Members:} Ali Ekin Özçetin, Ali Sezer Uysal, Mert Güner

\textbf{Problem Definition:} Large Language Models (LLMs) frequently generate code containing critical security vulnerabilities (CWEs). This project evaluates whether prompting strategy (zero-shot, few-shot, chain-of-thought) and model choice influence the frequency and resolvability of security issues in LLM-generated Python code.

\begin{itemize}[noitemsep]
    \item \textbf{Task:} Code generation paired with automated evaluation and iterative repair.
    \item \textbf{Importance:} As AI-assisted coding becomes ubiquitous, ensuring the security of generated code without human intervention is a critical challenge.
    \item \textbf{Pipeline:} The system receives a natural language instruction and outputs a Python script. If vulnerabilities are detected via SAST, tool outputs are fed back as input for the next generation cycle.
    \item \textbf{Technical Challenge:} Model self-assessment calibration — models must not only generate functional code but also accurately judge their own patches without falling into infinite repair loops or hallucinating false vulnerabilities.
\end{itemize}

% ════════════════════════════════════════════
\section{Dataset / Benchmark Status}
% ════════════════════════════════════════════

We use the \textbf{LLMSecEval} benchmark (Tony et al., 2023): 150 Python prompts across 18 CWE categories drawn from the CWE Top 25. For this mid-phase, \textbf{15 prompts} were selected via stratified sampling (\texttt{seed=42}), covering 12 distinct CWE categories. The subset is saved in \texttt{data/mid\_phase\_prompts.json} and shared as a control variable across both experiment branches.

\begin{table}[h]
\centering
\caption{CWE distribution of the 15 mid-phase prompts}
\begin{tabular}{lll}
\toprule
\textbf{CWE ID} & \textbf{Vulnerability Type} & \textbf{Count} \\
\midrule
CWE-78  & OS Command Injection             & 2 \\
CWE-89  & SQL Injection                    & 2 \\
CWE-502 & Deserialization of Untrusted Data & 2 \\
CWE-22  & Path Traversal                   & 1 \\
CWE-79  & Cross-Site Scripting             & 1 \\
CWE-434 & Unrestricted File Upload         & 1 \\
CWE-306 & Missing Authentication           & 1 \\
CWE-200 & Info Exposure                    & 1 \\
CWE-20  & Improper Input Validation        & 1 \\
CWE-732 & Incorrect Permissions            & 1 \\
CWE-522 & Insufficiently Protected Credentials & 1 \\
CWE-798 & Hard-coded Credentials           & 1 \\
\bottomrule
\end{tabular}
\label{tab:cwe_dist}
\end{table}

\textbf{Known dataset issue:} CWE-78 prompts exhibit strong adversarial bias, forcing almost all models toward \texttt{os.system()} usage, creating near-inevitable SAST triggers regardless of prompting strategy.

% ════════════════════════════════════════════
\section{Literature Review Progress}
% ════════════════════════════════════════════

Five papers directly inform the methodology: Pearce et al.\ \cite{pearce2022asleep} show GitHub Copilot produces vulnerable code in \textasciitilde{}40\% of scenarios, motivating iterative repair pipelines. Tony et al.\ \cite{tony2023llmseceval} introduce LLMSecEval and report a 72\% Bandit hit rate on ChatGPT, which is the direct baseline for comparison. Nishida et al.\ \cite{nishida2024codeenhancer} propose the CodeEnhancer architecture; our project implements and extends their Stage 1 pipeline. Brown et al.\ \cite{brown2020gpt3} and Wei et al.\ \cite{wei2022cot} provide the theoretical foundation for our few-shot and chain-of-thought prompt formulations, respectively.

These papers collectively justify our technical direction: evaluating whether prompting strategies can act as a substitute for computationally expensive fine-tuning in security-aware code generation.

Planned additions for the final report: Kojima et al.\ zero-shot CoT (2022) and HumanEval-Security benchmarks.

% ════════════════════════════════════════════
\section{Baseline Models and Prompting Strategies}
% ════════════════════════════════════════════

Five open-source models were evaluated across two experiment branches, all running locally via Ollama to eliminate API non-determinism:

\begin{itemize}[noitemsep]
    \item \textbf{Code-specialized:} Qwen 2.5 Coder 7B (Branch B), DeepSeek-Coder 6.7B (Branch A)
    \item \textbf{General-purpose:} Llama 3.1 8B (both branches), Gemma 2 9B (Branch B), Mistral 7B (Branch A)
\end{itemize}

Three prompting strategies are implemented in \texttt{prompts/}: \textbf{Zero-shot} (system prompt only, no security instruction), \textbf{Few-shot} (3 secure examples covering CWE-78, CWE-89, CWE-502, following Brown et al.\ \cite{brown2020gpt3}), and \textbf{Chain-of-Thought} (4-step security reasoning template, following Wei et al.\ \cite{wei2022cot}).

Each branch runs a 3×3 design (3 models × 3 strategies) × 15 prompts = \textbf{135 generated files per branch}, for a combined total of \textbf{270 evaluated files}.

% ════════════════════════════════════════════
\section{Initial Experimental Results}
% ════════════════════════════════════════════

All files were validated through the full pipeline: Bandit + Pylint SAST $\rightarrow$ LLM repair loop (MAX\_ATTEMPTS=5) $\rightarrow$ functional correctness check.

\begin{table}[h]
\centering
\caption{Aggregated results by model across both branches (averaged over 3 strategies, 15 prompts each)}
\begin{tabular}{llcccc}
\toprule
\textbf{Model} & \textbf{Type} & \textbf{Bandit@0 (\%)} & \textbf{Resolution (\%)} & \textbf{Unresolved (\%)} & \textbf{Avg Iter} \\
\midrule
Qwen 2.5 Coder 7B   & Code    & 55.6 & \textbf{88.9} & 11.1 & \textbf{1.91} \\
Mistral 7B          & General & 31.1 & 75.6          & 24.4 & 2.62          \\
Gemma 2 9B          & General & 26.7 & 66.7          & 33.3 & 3.04          \\
DeepSeek-Coder 6.7B & Code    & 28.9 & 20.0          & 80.0 & 4.53          \\
Llama 3.1 8B        & General & 44.4 & 15.6          & 84.4 & 4.73          \\
\bottomrule
\end{tabular}
\label{tab:model_summary}
\end{table}

\begin{table}[h]
\centering
\caption{Strategy-level averages across all 5 models (15 prompts each)}
\begin{tabular}{lcccc}
\toprule
\textbf{Strategy} & \textbf{Bandit@0 (\%)} & \textbf{Resolution (\%)} & \textbf{Unresolved (\%)} & \textbf{Avg Iter} \\
\midrule
Zero-shot        & 42.7 & 52.0 & 48.0 & 3.47 \\
Few-shot         & 37.3 & 54.7 & 45.3 & 3.28 \\
Chain-of-Thought & 32.0 & 53.3 & 46.7 & 3.36 \\
\bottomrule
\end{tabular}
\label{tab:strategy_summary}
\end{table}

\textbf{Key findings:}
\begin{enumerate}[noitemsep]
    \item \textbf{Code-specialized models lead in resolution} — Qwen 2.5 Coder (88.9\%) outperforms all others, confirming that code-domain pretraining improves iterative security repair. Notably, DeepSeek-Coder (20.0\%) does not share this advantage, suggesting that code specialization alone is insufficient without strong instruction-following.
    \item \textbf{Mistral 7B is the top general-purpose model} — 75.6\% resolution with the lowest average iteration count (2.62) among general-purpose models, suggesting strong instruction compliance in the repair loop.
    \item \textbf{CoT reduces initial vulnerability rate} — Chain-of-Thought prompting achieves the lowest Bandit@0 in Branch B (28.9\%), but resolution rates across strategies are comparable (35.5–37.8\% in Branch A; 51.1–53.3\% in Branch B), confirming that model capability dominates over prompting strategy in the repair phase.
    \item \textbf{CWE-78 is universally persistent} — OS Command Injection is the most frequently flagged category across all models and strategies, due to prompt-induced bias toward \texttt{os.system()} usage.
\end{enumerate}

% ════════════════════════════════════════════
\section{Planned Improvements and Technical Direction}
% ════════════════════════════════════════════

\begin{itemize}[noitemsep]
    \item \textbf{Decoupled judge architecture} — Using the generating model as its own judge introduces calibration bias. We plan to designate a single high-performing model (e.g., Qwen 2.5) as the universal LLM-Judge for all models.
    \item \textbf{Scale to all 150 prompts} — Expand experiments to the full LLMSecEval Python subset for statistical power.
    \item \textbf{Functional correctness} — Introduce pass@k via HumanEval to measure the security/functionality trade-off.
    \item \textbf{Expanded few-shot coverage} — Add CWE-434 and CWE-22 examples to reduce hit rates for those persistent categories.
\end{itemize}

% ════════════════════════════════════════════
\section{Ablation and Prompt Sensitivity Plan}
% ════════════════════════════════════════════

\begin{itemize}[noitemsep]
    \item \textbf{CoT depth} — 2-step vs.\ 4-step CoT on 15-prompt subset to measure cognitive load impact
    \item \textbf{Few-shot example order} — 3 permutations to measure sequence bias
    \item \textbf{Iteration cap} — MAX\_ATTEMPTS=3 vs.\ 5 to assess compute-to-resolution ratio
    \item \textbf{Cross-model judge} — Compare self-judge vs.\ Qwen-as-judge resolution rates for Llama~3.1 and DeepSeek-Coder to isolate calibration bias
\end{itemize}

% ════════════════════════════════════════════
\section{Error Analysis and Generation Quality Analysis}
% ════════════════════════════════════════════

Three recurring failure patterns were identified in files unresolved after 5 iterations:

\begin{enumerate}[noitemsep]
    \item \textbf{Loop divergence (judge calibration failure)} — The model passes SAST but the LLM-Judge returns ``Incorrect'' repeatedly. Observed in Llama 3.1 8B and DeepSeek-Coder; the judge demands arbitrary library usage (e.g., Flask-WTF) rather than accepting functionally correct fixes.
    \item \textbf{Partial / oscillating fix} — The model removes one Bandit issue but introduces another. Observed in DeepSeek-Coder for CWE-732/CWE-200 co-occurring files.
    \item \textbf{Strategy paradox (Gemma 2)} — Gemma 2's resolution rate \emph{decreases} under CoT (60.0\%) vs.\ Zero-shot (66.7\%), despite CoT reducing initial Bandit findings. Multi-step reasoning appears to degrade coherence during the repair phase for this model.
\end{enumerate}

A full qualitative analysis ($\geq$5 unresolved files per model per CWE) will be completed for the final report.

% ════════════════════════════════════════════
\section{Ethical Considerations and Bias}
% ════════════════════════════════════════════

\textbf{Dual-use:} All generated code is stored locally and used solely to measure vulnerability rates. LLMSecEval is a published academic benchmark; no novel exploit code is generated.

\textbf{Measurement bias:} Bandit is rule-based and produces false positives (e.g., \texttt{debug=True} flagged HIGH in development contexts) and false negatives (logic flaws). The resolution rate measures \emph{Bandit-compliance}, not true security.

\textbf{Model bias:} Performance differences may partly reflect English instruction-following quality rather than intrinsic security reasoning. Models differ in training data, instruction-tuning, and language mix.

% ════════════════════════════════════════════
\section{GitHub and Reproducibility Status}
% ════════════════════════════════════════════

Repository: \href{https://github.com/Aliekinozcetin/CodeEnhancer}{github.com/Aliekinozcetin/CodeEnhancer}. Branch A results are on \texttt{main}; Branch B results are on \texttt{AliSezer\_Experiments}. Both branches share \texttt{data/mid\_phase\_prompts.json}, all prompt templates, and \texttt{analysis/compare\_results.py} metric formulas as control variables.

To reproduce Branch A:
\begin{verbatim}
git clone https://github.com/Aliekinozcetin/CodeEnhancer && cd CodeEnhancer
python3 -m venv .venv && source .venv/bin/activate && pip install -r requirements.txt
ollama pull llama3.1 && ollama pull deepseek-coder:6.7b && ollama pull mistral:7b
python3 code_generator.py && python3 code_validator.py
python3 analysis/compare_results.py && python3 analysis/visualize.py
\end{verbatim}

% ════════════════════════════════════════════
\section{Current Challenges and Next Steps}
% ════════════════════════════════════════════

\textbf{Challenges:} (1) LLM-Judge calibration failure — self-assessment loop is unreliable for general-purpose models, inflating unresolved rates. (2) Compute time — 150 prompts × 5 models × 5 iterations scales to hundreds of hours of sequential local compute; batched overnight runs are planned. (3) CWE-200 false positives from logging statements inflate unresolved rates.

\textbf{Next steps:} Merge branch data logs; decouple LLM-Judge role; run full 150-prompt batch; conduct qualitative error analysis ($\geq$5 files per model); draft final report in LNCS LaTeX template.

% ════════════════════════════════════════════
\begin{thebibliography}{9}

\bibitem{pearce2022asleep}
H.~Pearce, B.~Ahmad, B.~Tan, B.~Dolan-Gavitt, R.~Karri,
``Asleep at the Keyboard? Assessing the Security of GitHub Copilot's Code Contributions,''
\textit{IEEE S\&P}, 2022.

\bibitem{tony2023llmseceval}
C.~Tony, M.~Mutas, N.~Díaz Ferreyra, R.~Scandariato,
``LLMSecEval: A Dataset of Natural Language Prompts for Security Evaluations,''
\textit{MSR}, 2023.

\bibitem{nishida2024codeenhancer}
K.~Nishida, T.~Shibata, H.~Iida,
``CodeEnhancer: Enhancing LLM-Generated Code for Security and Reliability,''
\textit{JAIST Technical Report}, 2024.

\bibitem{brown2020gpt3}
T.~Brown et al.,
``Language Models are Few-Shot Learners,''
\textit{NeurIPS}, 2020.

\bibitem{wei2022cot}
J.~Wei et al.,
``Chain-of-Thought Prompting Elicits Reasoning in Large Language Models,''
\textit{NeurIPS}, 2022.

\end{thebibliography}

\end{document}
